\begin{abstract}
{\color{red}%
\DIFdel{%
Deep learning (DL) frameworks rely heavily on operator libraries to support diverse model architectures and dynamic input shapes. To maintain flexibility, these libraries employ generic kernels with numerous runtime parameters for optimal scheduling. However, relying on runtime resolution for these parameters hurts performance, whereas statically enumerating them for efficiency leads to severe code bloat, creating a fundamental dilemma: accepting performance degradation with generic kernels or suffering severe code bloat through exhaustive specialization.
}
\par}
{\color{blue}%
Deep-learning (DL) systems rely on expert-written operator libraries for high performance across dynamic input shapes. These libraries use shape-dependent schedule parameters, such as tile sizes and on-chip buffer layouts. Schedule-generic expert-written operator libraries preserve shape coverage with a single kernel but limit compile-time optimization, whereas schedule-fixed expert-written operator libraries precompile combinations of schedule parameters and incur binary bloat. Kernel-generation compilers do not resolve this trade-off because they generate new kernels rather than specialize existing libraries.
\par}

{\color{red}%
\DIFdel{%
A promising solution is to specialize kernels by propagating model-specific constants from the framework to the compiler. However, this approach faces two challenges: (1) Cross-language semantic gap, where constant information in the high-level Python DSL fails to propagate to device-native C++ compilers, and (2) Control flow obstruction, where extensive validity checks on dynamic dimensions generate massive execution paths handling invalid shapes, thereby impeding the precision and efficiency of constant propagation.
}
\par}
{\color{blue}%
Although many schedule parameters are fixed for a given model, the corresponding model-specific constants recorded in Graph IR are absent from the LLVM IR of separately compiled C++ operator libraries. Runtime argument transmission through Python/C++ FFI does not make these values available as compile-time facts. Propagating model-specific constants from the framework enables specialization, but existing toolchains face two challenges: the Python/C++ semantic gap prevents constant information from reaching the native compiler, and shape checkers for dynamic dimensions introduce many early-exit paths that weaken constant propagation.
\par}

{\color{red}%
\DIFdel{%
We present TilingInfer, the first framework designed to automate kernel specialization in cross-language and dynamic-shape scenarios. To bridge the semantic gap, TilingInfer employs program synthesis to generate partial evaluation programs directly from the Graph IR. To handle control flow, it constructs a Conditional Return-Value Flow Graph to trace the relationship between shape check conditions and early-exit paths, enabling efficient identification and pruning of infeasible execution paths through lightweight constant comparison.
}
\par}
{\color{blue}%
We present TilingInfer, a static-analysis framework that uses LLVM to propagate Graph IR shape information through existing C++ operator libraries and derive constant schedule parameters for kernel specialization and redundant-code elimination. To address the two challenges above, TilingInfer introduces two core designs. First, type-specific context-update rules construct C++ entry states directly from Graph IR shapes and attributes, avoiding analysis of Python/C++ FFI glue code. Second, a Conditional Return-Value Flow Graph identifies early-exit paths and prevents their states from affecting normal paths, preserving precision without analyzing every path.
\par}

{\color{red}%
\DIFdel{%
We implement TilingInfer on Ascend NPUs and NVIDIA GPUs. Evaluations on state-of-the-art libraries (CANN and FlashInfer) demonstrate its effectiveness: TilingInfer achieves an average 1.06$\times$ speedup (up to 1.17$\times$) on Ascend operators and 1.03$\times$ speedup on LLM inference. Furthermore, it reduces the binary size of FlashInfer by over 96\% on GPUs, demonstrating significant potential for library slimming.
}
\par}
{\color{blue}%
On an Ascend NPU, Tiling achieves an average $1.06\times$ operator speedup across diverse operators and an average $1.05\times$ speedup across six representative inference scenarios. It provides larger benefits in the LLM decode phase and recommendation models, averaging a $1.10\times$ speedup. On an NVIDIA GPU, it reduces a 542\,MB binary size by 88.8--95.5\% across 17 model variants.
\par}
\end{abstract}
